\section{Introduction}
\label{sec:intro}
We are developing a system to make academic schedules for a semester at the University of Calgary, a single weekly schedule (which is then repeated for 13 weeks) containing $\mathit{Lectures}$ and $\mathit{Tutorials}$ sections (belonging to a course) which fit in weekly room time $\mathit{Slots}$ with designated start and end times. There are two categories of room time $\mathit{Slots}$, one for lectures and the other for tutorials, along with specific ``active learning'' designation for some rooms, with the number of rooms limited per slot. The objective is to assign lectures and tutorials to permitted weekly slots so that every hard constraint holds and the soft-constraint penalty is minimized. This problem is similar to what a student would do when creating a weekly class schedule in Schedule Builder, although we are creating the section times which students pick from. There is no direct step-based solution to this problem; therefore a procedural approach would require brute force. The number of possible timetables grows exponentially, about $|\mathit{Slots}|^{|X|}$. Search scales better in practice as it abandons a partial timetable as soon as it breaks a hard constraint.

This proposal compares two ways to improve a valid timetable: OR-tree feasibility followed by set-based local search, or OR-tree feasibility followed by AND-tree branch-and-bound. Section~\ref{sec:problem} formalizes the scheduling instance once; later sections instantiate the search models and controls without redefining the same constraints.